% =====================================================
% COMMON COMPLEXITY NOTATION
% =====================================================
\newcommand{\NP}{\mathrm{NP}}
\newcommand{\PClass}{\mathrm{P}}
\newcommand{\SAT}{\mathrm{SAT}}
\newcommand{\VC}{\mathrm{VC}}
\newcommand{\BTD}{\mathrm{BTD}}
\newcommand{\Z}{\mathbb{Z}}
\newcommand{\R}{\mathbb{R}}
\newcommand{\rT}{\sqrt{3}}
% \newcommand{\root3o2}{\frac{\sqrt{3}}{2}}

% =====================================================
% TITLE
% =====================================================
\title{Bloons Tower Defense is NP-Hard}
\author{farms4life2016}
\date{\today}

\setlength{\parindent}{0pt}

\begin{document}

\maketitle

\begin{abstract}
We prove that a decision version of generalized Bloons Tower Defense is NP-hard via a polynomial-time reduction from Monotone Rectilinear Planar 3-SAT.
\end{abstract}

\tableofcontents

\newpage

\section{Introduction}


\section{Definitions}

Lines can be described as equations, e.g. $L \coloneqq y = mx + b$.
Lines can also be described as a set of points, e.g. $L = \{ (x,y) \in \R^2 : y=mx+b \}$.
If we think of lines as sets, then we can union two lines to get a plane with both lines drawn on it, and we can intersect two lines to get the points where both lines intersect.

\subsection{Grid}

We will define a grid $G$ as a set of lines drawn on the Cartesian Plane.

\subsection{Square Grid}

Define $L_h$ as the union of all horizontal lines $y=c_1$ such that $(x,y) \in \R^2$ and $c_1 \in \Z$:
$$L_h = \{ (x, y) \in \R^2 : \forall c_1 \in \Z, y = c_1 \}$$

Define $L_v$ as the union of all vertical lines $y=c_2$ such that $(x,y) \in \R^2$ and $c_2 \in \Z$:
$$L_v = \{ (x, y) \in \R^2 : \forall c_2 \in \Z, x = c_2 \}$$

Define the \textit{square grid} $G_{\square}$ as the following union:
$$G_{\square} = L_h \cup L_v$$

Define the \textit{square grid intersections} $I_{\square}$ as the following intersection:
$$I_{\square} = L_h \cap L_v$$

\subsection{Square Grid Aligned}

A graph $G=(V,E)$ is \textit{square grid aligned} if it can be drawn on the Cartesian Plane with the following restrictions:
\begin{itemize}
    \item $G$ is planar
    \item for all $v \in V$ where $v = (v_x, v_y)$, $(v_x, v_y) \in I_{\square}$.
    \item for all $e \in E$ that connects two vertices $v_s, v_t \in V$, $e$ can be written as an ordered collection of line segments of size $p-1$ of the form:
    $$e = \{ \left( (x_1, y_1), (x_2, y_2) \right), \left( (x_2, y_2), (x_3, y_3) \right), \dots, \left( (x_{p-1}, y_{p-1}), (x_p, y_p) \right) \}$$
    In other words, $e$ is a list of pairs of points, with each pair defining the beginning and end of a line segment:
    \begin{itemize}[label=$\circ$]
        \item for all $i=1,2,\dots, p$, $(x_i, y_i) \in I_{\square}$.
        \item for all $j=2,3,\dots, p-1$, $(x_j, y_j) \notin V$.
        \item $v_s = (x_1, y_1)$.
        \item $v_t = (x_p, y_p)$.
        \item for all $k=1, 2, \dots, p-2$, the second point in the $k$-th pair is the same as the first point in the $(k+1)$-th pair.
        \item \textbf{Additionally, all line segments must have length 1}, i.e. \newline for all $q=1,2,\dots, p-1$, $(x_q - x_{q+1})^2 + (y_q - y_{q+1})^2 = 1$
    \end{itemize}
    
\end{itemize}

\subsection{Triangle Grid}

This definition is similar to its square counterpart, except we use different sets of lines. \newline

Define $L_s$ as the union of all horizontal lines $y = \frac{\rT}{2} c_1$ such that $(x,y) \in \R^2$ and $c_1 \in \Z$:
$$L_s = \{ (x, y) \in \R^2 : \forall c_1 \in \Z, y = c_1 \}$$
The $s$ subscript refers to the scaling factor of $ \frac{\rT}{2}$ compared to $L_h$.

Define $L_f$ as the union of all diagonal lines $y = \rT (x-c_2)$ such that $(x,y) \in \R^2$ and $c_2 \in \Z$:
$$L_f = \{ (x, y) \in \R^2 : \forall c_2 \in \Z, y = \rT (x-c_2) \}$$
The $f$ subscript stands for forward-slash (\verb|/|).

Define $L_b$ as the union of all diagonal lines $y = - \rT (x-c_3)$ such that $(x,y) \in \R^2$ and $c_3 \in \Z$:
$$L_b = \{ (x, y) \in \R^2 : \forall c_3 \in \Z, y = - \rT (x-c_3) \}$$
The $f$ subscript stands for back-slash (\verb|\|).

Define the \textit{triangle grid} $G_{\triangle}$ as the following union:
$$G_{\triangle} = L_s \cup L_f \cup L_b$$

Define the \textit{triangle grid intersections} $I_{\triangle}$ as the following intersection:
$$I_{\triangle} = L_s \cap L_f \cap L_b$$

\subsection{Triangle Grid Aligned}

This definition is similar to its square counterpart, except we use the triangle-related sets instead of the square-related ones.
\newline

A graph $G=(V,E)$ is \textit{triangle grid aligned} if it can be drawn on the Cartesian Plane with the following restrictions:
\begin{itemize}
    \item $G$ is planar
    \item for all $v \in V$ where $v = (v_x, v_y)$, $(v_x, v_y) \in I_{\triangle}$.
    \item for all $e \in E$ that connects two vertices $v_s, v_t \in V$, $e$ can be written as an ordered collection of line segments of size $p-1$ of the form:
    $$e = \{ \left( (x_1, y_1), (x_2, y_2) \right), \left( (x_2, y_2), (x_3, y_3) \right), \dots, \left( (x_{p-1}, y_{p-1}), (x_p, y_p) \right) \}$$
    In other words, $e$ is a list of pairs of points, with each pair defining the beginning and end of a line segment:
    \begin{itemize}[label=$\circ$]
        \item for all $i=1,2,\dots, p$, $(x_i, y_i) \in I_{\triangle}$.
        \item for all $j=2,3,\dots, p-1$, $(x_j, y_j) \notin V$.
        \item $v_s = (x_1, y_1)$.
        \item $v_t = (x_p, y_p)$.
        \item for all $k=1, 2, \dots, p-2$, the second point in the $k$-th pair is the same as the first point in the $(k+1)$-th pair.
        \item \textbf{Additionally, all line segments must have length 1}, i.e. \newline for all $q=1,2,\dots, p-1$, $(x_q - x_{q+1})^2 + (y_q - y_{q+1})^2 = 1$
    \end{itemize}

\subsection{Triangle Grid Embedded}

A graph $G=(V,E)$ is \textit{triangle grid embedded} if it is triangle-aligned and $p=2$, meaning every $e \in E$ consists of a single line segment of length 1 directly connecting two vertices.
    
\end{itemize}

\section{Data Structures}

To avoid the additional computation and inaccuracies associated with floating point numbers,
we want to store all numbers as integers. We note that in the aforementioned definitions, $G=(V,E)$ can be stored entirely using some collection(s) of points in $I_\square$ or some collection(s) of points in $I_\triangle$ (but not both).
Thus, we need to find a way to represent each point in $I_\square$ and $I_\triangle$ using only integers.

For $I_\square$, this is easy. Every $(x,y) \in I_{\square}$ satisfies $(x,y) \in \Z^2$, so we can use $\Z^2$ to represent points in $I_\square$.

For $I_\triangle$, we can still use the $\Z^2$. We assume that the solver or verifier does not care about the drawing of the graph, only the data. (This can probably be formally proved using the equivalence/isomorphism of graphs... todo later...).
From $\Z^2$, we can transform into $I_\triangle$ by first adding $\frac{1}{2}$ to each odd x-coordinate. Then, scale all $y$ coordinates by a factor of $\frac{\rT}{2}$. (TODO: prove why this works, probably using a bijection)

There are probably many other ways of doing this...some of which are easier to prove.

Therefore, we can use $\Z^2$ to represent $G$ in both square and triangle cases.

\section{Decision Problems}

\subsection{Monotone (Rectilinear) Planar 3SAT (MRP3SAT)}

This is a known NP-Complete problem.

Any Planar 3SAT formula has a \textit{rectilinear embedding}, which means you can draw the incidence graph such that
variables are on a \textit{variable row} as horizontal lines,
clauses are above or below the variable row as horizontal lines,
and vertical lines connect clauses and variables if and only if a variable or its negation appears in that clause.
(This is also a known result.)

Monotone just means that all clauses are either positive or negative, 
meaning that the clause contains all non-negated variables or all negated variables, respectively.
Additionally, all positive clauses must be drawn above the variable row, and all negative clauses must be drawn below the variable row.

Importantly, clauses can contain duplicate literals and can only contain at most 3 literals.
(If every clause only has 3 distinct variables, then the answer is always ``yes'' and we can find a certificate in polytime.
This is also a known result.)

\subsection{Square Grid Aligned MRP3SAT (SGA MRP3SAT)}

MRP3SAT but the incidence graph $G$ is square grid aligned.

\subsection{Triangle Grid Aligned MRP3SAT 3SAT (TGA MRP3SAT)}

MRP3SAT but the incidence graph $G$ is triangle grid aligned.

\subsection{Triangle Grid Aligned Planar Vertex Cover (TGA PVC)}

Planar Vertex Cover but the graph $G$ is triangle grid aligned.
Our reduction would still work if the max. in-degree is 3 (verify).

\subsection{Triangle Grid Embedded Planar Vertex Cover (TGE PVC)}

Planar Vertex Cover but the graph $G$ is triangle grid embedded.
Our reduction would still work if the max. in-degree is 3.
